% =====================================================================
%  Chapter 2 — How We Got Here
%  Map: research/chapter-map.md (Chapter 2); main source research/silica-history.md
%  Rule: the author is never quoted in body text; short quotes only in
%  footnotes. Full statements live in interviews/.
%  Caution: "Phase 1/2" of the compiler ≠ "Phases 2–8" of chunk 15.
% =====================================================================
\chapter{How We Got Here}
\label{ch:history}

Silica's repository holds 686 commits made between 26~December 2025 and
13~November 2026. A year is not long for a programming language. But the ideas
in those commits are older than the repository, and the story starts earlier,
in four projects that each tried to fix something from the inside.

\section{Before Silica}

\subsection{SwErl, 2023--2025}

SwErl began in February 2023 as a library that brought Erlang's model of
concurrency into Swift: lightweight processes, message sending, and the
standard design patterns Erlang programmers rely on. Its argument was that async/await and promises try to make
concurrent code look sequential, and that this works until something goes
wrong.\footnote{SwErl README, ``Painless Concurrency'':
\src{https://github.com/inter-erlang/SwErl}.} Students contributed to it, and in
2024 a branch could already speak Erlang's distribution protocol.

SwErl also showed what a library cannot do. Swift's own dictionaries crashed
under concurrent access, so the registry of processes had to be rebuilt more
than once. There were no supervisors. Messages were untyped, so every handler
began by casting. The networking code only compiled on Apple platforms. In
September 2025 the unfinished node work was handed to a different project.

% TODO-INTERVIEW: What did SwErl teach you that you could only have learned by
%   building it inside someone else's language? Is Silica SwErl's answer?

\subsection{Unsafe Impedance, 2024}

The paper described in Chapter~\ref{ch:why} began with a request from a member
of the Erlang Ecosystem Foundation at CodeBEAM in San Francisco. It was outlined
in March 2024, submitted in May and presented in Milan that September. Its first
outline celebrated how Erlang and Elixir avoid common memory errors. The finished
paper followed the evidence further: native functions loaded into the BEAM can
break every guarantee the VM offers, and they can be hidden from the Erlang code
that calls them.

% TODO-VERIFY: CodeBEAM San Francisco year (likely March 2024).

\subsection{BEAM\_SEMP, 2025}

From late August to October 2025, Barney built a secure overlay for Erlang
networking. Its trust layer used mutual TLS, admitted peers only by certificate
fingerprint, and kept a deny list that overrode even trusted peers: no remote
code loading, no atom creation, no loading of native code. One module switched
off Erlang's standard distribution entirely. Errors were never sent back to a
peer. It worked, and every part of it was a patch on top of permissive defaults.

\subsection{A Rust BEAM, and a recovery, 2025}

Through 2025, Barney was also porting the BEAM's virtual machine from C to Rust,
to answer the people who asked how Erlang could be memory safe when its virtual
machine was written in C. In June of that year Barney had a mini-stroke, which
for a time took away the ability to read, to think about complex things and to
form sentences easily. The work on the port became part of the recovery, and
language models wrote much of the code Barney could not yet write. On
24~December 2025, well into the port, Barney realised it was the wrong problem
to be solving. Chapter~\ref{ch:why} tells that story.

\section{Scaffolding}

The first commit, on 26~December 2025, was a bootstrap compiler for AArch64
written in Rust that emitted LLVM code. The founding notes already set the tone:
no C interoperability, no loops, recursion only, and a syntax meant to be
robust for language models to generate. Within a day the LLVM intermediate
files were gone from the repository; within six days, so were generics.

LLVM was a scaffold from the start. The plan was always for the real compiler
to emit machine code itself, so that it could reason about the facts of the
chip it was targeting (Chapter~\ref{ch:machines}).

% TODO-INTERVIEW: Why Rust for the bootstrap, given later views on Rust?

\section{A language writing itself}

On 17~January 2026 work began on the second compiler, written in Silica and
built by the Rust one. On 16~February it produced its first working program:
a trivial program that returned the number 42, compiled straight to machine instructions with no LLVM in the
path. The test suite, called the \emph{trials}, started three days later. It
would become the project's arbiter of truth (Chapter~\ref{ch:refuses}).

Spring brought the language's core: regions in March, sequence blocks for
effects, actors in April, then modules, traits and supervisors. In mid-April the
project went public with a README, a licence and a motto. An AI policy followed
at the end of the month. In May came the foreign function interface, Fifi,
reversing the founding refusal of C interop---but wrapped in so much visible
danger that the reversal became one of the language's principles
(Chapter~\ref{ch:danger}).

On 29~June the entire first data-structure library was deleted as too
complicated for its purpose, and a smaller, trait-based design replaced it.

\section{Standing on its own}

Through July 2026 the compiler was converted to build itself, a process
punctuated by machines running out of memory. At the end of July the first
fully self-hosted compiler binary compiled itself. On 12~September the compiler
reached a \emph{fixed point}: rebuilt by itself, it produced itself byte for
byte, while passing more than 33,000 trials. Five days later the same was true
on a Raspberry Pi, and the Rust bootstrap was deleted.

% TODO-INTERVIEW: Take me to the day of the first fixed point.

\section{Going everywhere}

Linux on x86-64 reached its own fixed point within days. The same week, work
began on a bare-metal port to the ESP32-S3 microcontroller, without the vendor's
SDK or operating system. In October, more than twenty-two thousand trials were
rewritten into supervisor-and-actor form; a thousand actors ran on a board with
512~kilobytes of memory; and on hosted systems, actors moved off one operating
system thread each onto shared carrier threads. On 19~October a single machine
held a million idle actors in 630 megabytes.

\section{The reckoning}

In late October, designing a revision of the standard data structures, the
project found that the runtime had been using a per-actor heap that no one had
ever decided on---a leftover of the first actor runtime---while the documentation
promised memory on actor stacks. Some forty design decisions were made in three
days. Everything moved onto actor stacks, with release proven by the compiler
rather than guessed at runtime (Chapter~\ref{ch:memory}). The new model was
built across all hosted targets by 4~November.

% TODO-VERIFY: Chunk 15 origin quotes are agent-recorded (.claude-memory);
%   confirm the account with the author before using any wording.

\section{Reaching out}

In the first two weeks of November, transport became part of the language
(Chapter~\ref{ch:transport}). Silica nodes spoke Erlang's wire format; boards
talked over serial lines; a Mac and a Linux machine talked over Ethernet with
mutual TLS. New targets appeared: an STM32 Nucleo board and a Raspberry Pi Pico
2~W. Alongside Silica, Project BEES---begun in April 2026 to let Erlang, Elixir,
Gleam and LFE be compiled to Silica---was pushing the language toward the
distribution story SwErl had started years before.

\section{How the work was done}

Three features of the process run through the whole record. The trials were the
gate: a change was not finished until every platform it could affect passed,
and reference assembly was never regenerated automatically. Barney made the
design decisions, many of them logged by date in the project's documents. And
from March 2026 onward, much of the implementation was done with AI coding
agents, assigned by task, working under rules written for them, one build at a
time because of memory. Barney thinks of them as an army of developers working
under direction, not as a way to improvise code, and trust in them grew as the
models improved and as more Silica existed for them to learn from. After the
October 2026 discovery that the army had built something no one had decided on,
the reviews multiplied: automated where possible, and repeated against the
design documents (Chapter~\ref{ch:hidden}).
% TODO-INTERVIEW: The quiet summer (June–August 2026): what was happening?
